\documentclass [aspectratio=169]{beamer}
\usetheme{Boadilla}
\usepackage{textpos} % package for the positioning
\usepackage[]{graphicx}
\usepackage{graphicx}
\usepackage{float}
\usepackage{hyperref}
\usepackage{caption}
\usepackage{subcaption}
\usepackage{algorithm,algpseudocode}
\usepackage{amsmath, amsfonts, amsthm, amssymb}
\usepackage{dsfont}
\usepackage[export]{adjustbox}
\usepackage{tikz}
\usetikzlibrary{positioning}
\usetikzlibrary{arrows, shapes, decorations, automata, backgrounds, fit, petri, calc}
\setbeamertemplate{itemize items}[circle]
\setbeamertemplate{enumerate items}[circle]

\newcommand*{\logofont}{\fontfamily{phv}\selectfont}
\definecolor{uoftblue}{RGB}{0,42,92} % official blue color for uoft
\definecolor{deptgreen}{RGB}{114,192,148} 
\definecolor{deptoran}{RGB}{252,103,63} 

\hypersetup{
  colorlinks   = true, %Colours links instead of boxes
  urlcolor     = uoftblue, %Colour for external hyperlinks
  linkcolor    = black, %Colour of internal links
  citecolor   = black %Colour of citations
}

% lin alg
\newcommand{\bu}{{\mathbf{u}}}
\newcommand{\bv}{{\mathbf{v}}}
\newcommand{\bw}{{\mathbf{w}}}
\newcommand{\bx}{{\mathbf{x}}}
\newcommand{\by}{{\mathbf{y}}}
\newcommand{\bz}{{\mathbf{z}}}
\newcommand{\zerovec}{{\mathbf{0}}}
\newcommand{\innerprod}[1]{\langle #1 \rangle}
\newcommand{\Tr}{\mathrm{Tr}}

% other useful stuff
\newcommand{\Id}{{\mathds{1}}}
\newcommand{\R}{{\mathbb{R}}}
\newcommand{\C}{{\mathbb{C}}}
\newcommand{\Z}{{\mathbb{Z}}}
\newcommand{\N}{{\mathbb{N}}}
\newcommand{\Q}{{\mathbb{Q}}}
\newcommand{\F}{{\mathbb{F}}}
\newcommand{\cL}{{\mathcal{L}}}
\newcommand{\cP}{\mathcal{P}}
\newcommand{\cT}{\mathcal{T}}
\newcommand{\inv}{{-1}}

% set theory
\newcommand{\interior}{\accentset{\circ}}



%\DeclareMathOperator{\dim}{dim}
\DeclareMathOperator{\range}{range}
\DeclareMathOperator{\rank}{rank}
\DeclareMathOperator{\nullspace}{null}
\DeclareMathOperator{\nullity}{nullity}

\beamertemplatenavigationsymbolsempty

% block
% example block
% alert block


\title[]{Module 9: Linear Algebra III \\ {\large Operational math bootcamp}\\ \includegraphics[width=8cm]{dept_logo.png}\vspace{-1em}}
\author[]{Emma Kroell}
\institute[]{University of Toronto}
\date{July 25, 2022}

% set color
\setbeamercolor{title in head/foot}{bg=white}
\setbeamercolor{author in head/foot}{bg=white}
\setbeamercolor{date in head/foot}{fg=uoftblue}
\setbeamercolor{date in head/foot}{bg=white}
\setbeamercolor{title}{fg=uoftblue}
\setbeamerfont{title}{series=\bfseries}
\setbeamercolor{frametitle}{fg=uoftblue}
\setbeamerfont{frametitle}{series=\bfseries}
\setbeamercolor*{item}{fg=uoftblue}
\setbeamercolor{block title}{bg=uoftblue}
\setbeamercolor{block title}{fg=white}
\setbeamercolor{block body}{bg=uoftblue!9!white}
\setbeamercolor{block title example}{bg=deptgreen}
\setbeamercolor{block title example}{fg=white}
\setbeamercolor{block body example}{bg=deptgreen!13!white}
\setbeamercolor{block title alerted}{bg=deptoran}
\setbeamercolor{block title alerted}{fg=white}
\setbeamercolor{block body alerted}{bg=deptoran!10!white}


% set logo at non-title pages
\logo{\includegraphics[height=0.8cm]{dept_logo.png}\vspace*{-.045\paperheight}\hspace*{.78\paperwidth}}

% set margin
\setbeamersize{text margin left=10mm,text margin right=10mm}

\begin{document}
{
\setbeamertemplate{logo}{}
\begin{frame}
    %\vspace{0.5in}
    \titlepage
    %\begin{textblock*}{10cm}(3.5cm,-7.5cm)
      %  \includegraphics[width=8cm]{dept_logo.png}
    %\end{textblock*}
\end{frame}
}

\begin{frame}{Outline}
Spectral theory and matrix decompositions:
\vspace{0.5em}
    \begin{itemize}
      \setlength\itemsep{0.5em}
    	\item Eigenvalues and eigenvectors
	\item Algebraic and geometric multiplicity of eigenvalues
	\item Jordan canonical form
	\item Singular value decomposition
	\item $LU$ and $QR$ decompositions
    \end{itemize}
\end{frame}

\begin{frame}{Recall: connection between matrices and linear maps}
\begin{block}{Multiplication by a matrix defines a linear map}
Let $A\in M_{m\times n}$ be a fixed matrix. Then, we can define a linear map $T_A \colon \F^n \to \F^m$ via $T_A(\bv) = A \bv$, where we recall matrix vector multiplication $(A\bv)_i = \sum_{k=1}^n A_{ik}v_k$ for $i=1, \ldots, m$.
\end{block}

\vspace{1em}

\begin{block}{Given a bases for $U$ and $V$, $T: U \to V$ can be written as a matrix}
Let $T \in \mathcal{L}(U,V)$ where $U$ and $V$ are vector spaces. Let $\bu_1, \ldots, \bu_n$ and $\bv_1, \ldots, \bv_m$ be bases for $U$ and $V$ respectively. The matrix of $T$ with respect to these bases is the $m \times n$ matrix $\mathcal{M}(T)$ with entries $A_{ij}$, $i = 1, \ldots, m$, $j = 1, \ldots, n$ defined by
\begin{equation*}
    T\bu_k = A_{1k} \bv_1 + \cdots + A_{mk} \bv_m.
\end{equation*}
\end{block}
\end{frame}

\begin{frame}{Eigenvalues}
\begin{definition}
Given an operator $A \colon V \to V$ and $\alpha \in \F$, $\lambda$ is called an \emph{eigenvalue} of $A$ if there exists a non-zero vector $\bv \in V\setminus\{\zerovec\}$ such that 
$$A \bv = \lambda \bv.$$
We call such $\bv$ an \emph{eigenvector} of $A$ with eigenvalue $\lambda$. We call the set of all eigenvalues of $A$ \emph{spectrum} of $T$ and denote it by $\sigma(T)$.
\end{definition}

\vspace{1em}

Motivation in terms of linear maps: Let $T \colon V \to V$ be a linear map, where $V$ is a vector space.  We would like to describe the action of this linear map in a particularly ``nice'' way: such that $T$ acts only by scaling, i.e. $T\bv_i = \lambda_i \bv_i$ where $\lambda_i \in \F$ for $i = 1, \ldots, n$.
\end{frame}

\begin{frame}{Finding eigenvalues}

\begin{itemize}
\item Rewrite $A \bv = \lambda \bv$ as  %$(A - \lambda I) \bv= 0$. 
\item Thus, if $\lambda$ is an eigenvalue, we can find the corresponding eigenvectors by finding the null space of $A - \lambda I$.
\item The subspace $\nullspace (A - \lambda I)$ is called the \emph{eigenspace}
\item To find the eigenvalues of $A$, one must find the scalars $\lambda$ such that $\nullspace (A - \lambda I)$ contains non-trivial vectors (i.e. not $\zerovec$)
\item Recall: We saw that $T \in \cL(U,v)$ is injective if and only if $\nullspace T = \{ \zerovec \}$.
\item Thus $\lambda$ is an eigenvector if and only if $A - \lambda I$ is not invertible.
\item Recall: $|A| \neq 0$ if and only if $A$ is invertible.
\item Thus $\lambda$ is an eigenvector if and only if %$|A - \lambda I|=0$.
\end{itemize}

\end{frame}


\begin{frame}
\begin{theorem}
The following are equivalent
\begin{enumerate}
    \item $\lambda\in \F$ is an eigenvalue of $A$,
    \item $(A-\lambda I)\bv = 0$ has a non-trivial solution,
    \item $|A-\lambda I|= 0$.
\end{enumerate}
\end{theorem}
\end{frame}


\begin{frame}{Characteristic polynomial}
\begin{definition}
If $A$ is an $n\times n$ matrix, $p_A(\lambda) = |A-\lambda I|$ is a polynomial of degree $n$ called the \emph{characteristic polynomial} of $A$.
\end{definition}

To find the eigenvectors of $A$, one needs to find the roots of the characteristic polynomial.
\end{frame}


\begin{frame}{Example}
Find the eigenvalues of
$$ \begin{bmatrix} 4 & -2 \\ 5 & -3 \end{bmatrix}.$$
\vspace{4cm}

\end{frame}


\begin{frame}{Multiplicity}
\begin{definition}
The multiplicity of the root $\lambda$ in the characteristic polynomial is called the \emph{algebraic multiplicity} of the eigenvalue $\lambda$. The dimension of the eigenspace $\nullspace (A - \lambda I)$ is called the \emph{geometric multiplicity} of the eigenvalue $\lambda$.
\end{definition}

\end{frame}


\begin{frame}
\begin{definition}[Similar matrices]
Square matrices $A$ and $B$ are called \emph{similar} if there exists an invertible matrix $S$ such that
$$A = SBS^\inv.$$
\end{definition}

\vspace{1em}

Similar matrices have the same characteristic polynomials and hence the same eigenvalues (see exercise).
\end{frame}


\begin{frame}
\begin{theorem}
Suppose $A$ is a square matrix with distinct eigenvalues $\lambda_1, \ldots, \lambda_n$. Let $\bv_1, \ldots, \bv_n$ be eigenvectors corresponding to these eigenvalues. Then $\bv_1, \ldots, \bv_n$ are linearly independent.
\end{theorem}
\begin{block}{Proof}
\vspace{3cm}
\end{block}
\end{frame}

\begin{frame}
\begin{block}{Proof continued}
\vspace{6cm}
\end{block}
\end{frame}


\begin{frame}
\begin{corollary}
If a $A\in M_n(\C)$ has $n$ distinct eigenvalues, then $A$ is diagonalizable. That is there exists an invertible matrix $S\in M_n(\C)$ such that $A = SDS^\inv$, where $D$ is a diagonal matrix with the eigenvalues of $A$ in the diagonal.
\end{corollary}

\end{frame}


\begin{frame}
\begin{theorem}
Let $A:V\to V$ be an operator with $n$ eigenvalues. $A$ is diagonalizable if and only if for each eigenvalue $\lambda$, the geometric multiplicity of $\lambda$ and the algebraic multiplicity of $\lambda$ are the same.
\end{theorem}
\end{frame}


\begin{frame}{Example: a diagonalizable matrix}
$$ \begin{bmatrix} 1 & 2 \\ 8 & 1\end{bmatrix} \; \text{is diagonalizable.}$$.
\vspace{4cm}
\end{frame}


\begin{frame}{Example continued}

\end{frame}


\begin{frame}{Example: a matrix that is not diagonalizable}
$$ \begin{bmatrix} 1 & 1 \\ 0 & 1\end{bmatrix} \; \text{is \emph{not} diagonalizable.}$$.
\vspace{4cm}
\end{frame}


\begin{frame}
\begin{theorem}
Let $A\in M_n(\R)$ be a symmetric matrix. Then, there exists an orthogonal matrix $O\in M_n(\R)$ such that $A=ODO^T$, where $D$ is a diagonal matrix with the eigenvalues of $A$ in the diagonal. Furthermore, all eigenvalues of $A$ are real.
\end{theorem}

\vspace{1em}

We can also state this for $M_n(\C)$: \\
Let $A\in M_n(\C)$ be a Hermitian matrix. Then, there exists a unitary matrix $U\in M_n(\C)$ such that $A = UDU^*$, where $D$ is a diagonal matrix with the eigenvalues of $A$ in the diagonal. Furthermore, all eigenvalues of $A$ are real.
\end{frame}


\begin{frame}{Block matrices}
\begin{definition}
A block matrix is a matrix that can be broken into sections called blocks, which are smaller matrices.
\end{definition}

\begin{example}
$$ \begin{bmatrix} 2 & 1 & 0 & 1  \\ 0 & 2 &1 & 1 \\ 0 & 0 & 1 & 0 \\ 0 & 0 & 2 & 1 \end{bmatrix} \qquad \qquad \qquad \qquad$$
\vspace{2cm}
\end{example}

\end{frame}


\begin{frame}

\begin{definition}
A square matrix is called \emph{block diagonal} if it can be written as a block matrix where the main-diagonal blocks are all square matrices and the off-diagonal blocks are all zero.
\end{definition}

\begin{example}
The matrix
$$ \begin{bmatrix} 2 & 1 & 0 & 0  \\ 0 & 2 &0 & 0 \\ 0 & 0 & 1 & 0 \\ 0 & 0 & 2 & 1 \end{bmatrix} $$
is block diagonal.
\end{example}
\end{frame}



\begin{frame}
\begin{definition}
A vector $\bv$ is called a \emph{generalized eigenvector} of $A$ corresponding to an eigenvalue $\lambda$ if there exists $k \geq 1$ such that 
$$ (A - \lambda I)^k \bv = 0.$$
\end{definition}
The set of generalized eigenvectors of an eigenvalue $\lambda$ (plus $\zerovec$) is called the \emph{generalized eigenspace} of $\lambda$.

\vspace{1em}

\begin{exampleblock}{Proposition}
The algebraic multiplicity of an eigenvalue $\lambda$ is the same as the dimension of the corresponding generalized eigenspace.
\end{exampleblock}

\end{frame}


\begin{frame}
\begin{theorem}[Jordan decomposition theorem]
For any operator $A$ there exists a basis such that $A$ is block diagonal with blocks that have eigenvalues on the diagonal and 1s on the upper off-diagonal. In other words, $A$ can be written in the form
$$ A = \begin{bmatrix} J_1 & 0 & \ldots & 0\\ 0 & J_2 & \ldots & 0 \\ \vdots & \vdots & \ddots & \vdots \\ 0 & 0 & \ldots & J_k \end{bmatrix} $$
where the blocks $J_i$ on the main diagonal are Jordan block of the form 
$$\begin{bmatrix} \lambda  \end{bmatrix}, \begin{bmatrix} \lambda & 1  \\ 0 &\lambda \end{bmatrix}, \begin{bmatrix} \lambda & 1 & 0 \\ 0 &\lambda & 1 \\ 0 & 0 & \lambda \end{bmatrix}, \text{ etc.}$$
This form is called \emph{Jordan canonical form}.
\end{theorem}
\end{frame}


\begin{frame}
Connection to algebraic and geometric multiplicity:
\begin{itemize}
    \item The algebraic multiplicity of an eigenvalue $\lambda$ is the number of times $\lambda$ appears on the diagonal.
    \item The geometric multiplicity of $\lambda$ is the number of Jordan blocks associated with $\lambda$.
\end{itemize}
\vspace{1em}

Why is Jordan form useful?

\vspace{4cm}

\end{frame}


\begin{frame}{Singular value decomposition}
\begin{itemize}
\item $A^T A$ is symmetric
\item therefore it is orthogonally diagonalizable and has real eigenvalues 
\item In fact, the eigenvalues are non-negative (exercise)
\end{itemize}


\begin{definition}
Let $A$ be an $m \times n$ matrix.  Let $\lambda_1, \ldots, \lambda_n$ be the eigenvalues of $A^T A$. Then the \emph{singular values} of $A$ are defined as
\begin{equation*}
    \sigma_1 = \sqrt{\lambda_1}, \ldots, \sigma_n = \sqrt{\lambda_n}.
\end{equation*}
\end{definition}
\end{frame}

\begin{frame}
\begin{theorem}[Singular value decomposition]
If $A$ is an $m \times n$ matrix of rank $k$, then we can write
$$ A = U \Sigma V^T$$
where $\Sigma$ is an $m \times n$ matrix of the form
$$ \begin{bmatrix} D_{k \times k} & 0_{k \times (n - k)} \\ 0_{(m -k) \times k} & 0_{(m -k) \times (n-k)} \end{bmatrix},$$
$D$ is a diagonal matrix with the singular values of $A$, $ \sigma_1, \ldots, \sigma_k$, on the diagonal and $U$ and $V$ are both orthogonal matrices (of size $m \times m$ and $n \times n$, respectively).
\end{theorem}
\end{frame}

\begin{frame}
Uses of SVD:

\vspace{3cm}

Differences between JCF and SVD:

\vspace{3cm}

\end{frame}

\begin{frame}{$LU$-decomposition}
\begin{definition}
The $LU$-decomposition of a square matrix $A$ is the factorization of $A$ into a lower triangular matrix $L$ and an upper triangular matrix $U$ as follows:
\begin{equation*}
    A = LU.
\end{equation*}
\end{definition}

Why is this useful? Consider the linear system $A \bx = \mathbf{b}$
\vspace{3cm}
\end{frame}

\begin{frame}{Recall: orthonormal basis}
\begin{definition}
Two vectors $\bx,\by \in V$ are called \emph{orthogonal} if $\innerprod{\bx,\by} = 0$, denoted $\bx \perp \by$. We call them \emph{orthonormal} if additionally the vectors are normalized, i.e. $\Vert \bx \Vert =\Vert \by\Vert = 1 $. A basis $\bx_1,\ldots, \bx_n$ of $V$ is called \emph{orthonormal basis (ONB)}, if the vectors are pairwise orthogonal and normalized.  
\end{definition}

\vspace{1em}

Starting from a set of linearly independent vectors, we can construct another set of vectors which are orthonormal and span the same space using the \textbf{\textcolor{deptoran}{Gram-Schmidt Algorithm}}.
\end{frame}



\begin{frame}{$QR$-decomposition}
\begin{definition}[$QR$-decomposition]
The $QR$-decomposition of an $m \times n$ matrix $A$ with linearly independent column vectors is the factorization of $A$ as follows:
\begin{equation*}
    A = QR,
\end{equation*}
where $Q$ is an $m \times n$ matrix with orthonormal column vectors and $R$ is an $n\times n$ invertible upper triangular matrix.
\end{definition}
\end{frame}

\begin{frame}
One obtains the factorization by applying the Gram-Schmidt algorithm to the columns of $A$. Let $\bu_1, \ldots, \bu_n$ be the column vectors of $A$. Let $\mathbf{q}_1, \ldots, \mathbf{q}_n$ be the orthonormal vectors obtained by applying Gram Schmidt. Then one can write:
\begin{align*}
    \bu_1 &= \innerprod{\bu_1,\mathbf{q}_1}\mathbf{q}_1 + \innerprod{\bu_2,\mathbf{q}_2}\mathbf{q}_2+ \ldots +  \innerprod{\bu_n,\mathbf{q}_n} \mathbf{q}_n \\
    \bu_2 &= \innerprod{\bu_2,\mathbf{q}_2}\mathbf{q}_2 + \ldots +  \innerprod{\bu_n,\mathbf{q}_n} \mathbf{q}_n \\
    & \vdots \\
    \bu_n &=  \innerprod{\bu_n,\mathbf{q}_n} \mathbf{q}_n
\end{align*}

Thus the orthonormal vectors obtained using Gram-Schmidt form the columns of $Q$, while $R$ is the terms needed to go between the columns of $A$ and thsoe of $Q$, i.e.
$$ R = \begin{bmatrix}
\innerprod{\bu_1,\mathbf{q}_1} & \innerprod{\bu_2,\mathbf{q}_2} & \ldots & \innerprod{\bu_n,\mathbf{q}_n} \\
0 & \innerprod{\bu_2,\mathbf{q}_2} & \ldots & \innerprod{\bu_n,\mathbf{q}_n} \\
\vdots &\vdots & \ddots & \vdots \\
0 & 0 & \ldots & \innerprod{\bu_n,\mathbf{q}_n} 
\end{bmatrix}.$$
\end{frame}

\begin{frame}
Why use $QR$-decomposition? 

\vspace{3cm}

\end{frame}

\begin{frame}{References}

Howard Anton and Chris Rorres. Elementary Linear Algebra. 11th ed. Wiley, 2014

\vspace{1em}

Axler S. \textit{Linear Algebra Done Right}. 3rd ed. Undergraduate Texts in Mathematics. Springer, 2015.
Available from: \href{https://link.springer.com/book/10.1007/978-3-319-11080-6}{https://link.springer.com/book/10.1007/978-3-319-11080-6} 
%U of T login:  \href{https://link-springer-com.myaccess.library.utoronto.ca/book/10.1007/978-3-319-11080-6}{https://link-springer-com.myaccess.library.utoronto.ca/book/10.1007/978-3-319-11080-6}

\vspace{1em}


\indent Treil S. \textit{Linear Algebra Done Wrong}. 2017. Available from: \href{https://www.math.brown.edu/streil/papers/LADW/LADW.html}{https://www.math.brown.edu/streil/papers/LADW/LADW.html}
\end{frame}




\end{document}